%%=============================================================================
%% Methodologie
%%=============================================================================

\chapter{\IfLanguageName{dutch}{Methodologie}{Methodology}}%
\label{ch:methodologie}

%% TODO: In dit hoofstuk geef je een korte toelichting over hoe je te werk bent
%% gegaan. Verdeel je onderzoek in grote fasen, en licht in elke fase toe wat
%% de doelstelling was, welke deliverables daar uit gekomen zijn, en welke
%% onderzoeksmethoden je daarbij toegepast hebt. Verantwoord waarom je
%% op deze manier te werk gegaan bent.
%% 
%% Voorbeelden van zulke fasen zijn: literatuurstudie, opstellen van een
%% requirementsanalyse, opstellen long-list (bij vergelijkende studie),
%% selectie van geschikte tools (bij vergelijkende studie, "short-list"),
%% opzetten testopstelling/PoC, uitvoeren testen en verzamelen
%% van resultaten, analyse van resultaten, ...
%%
%% !!!!! LET OP !!!!!
%%
%% Het is uitdrukkelijk NIET de bedoeling dat je het grootste deel van de corpus
%% van je bachelorproef in dit hoofstuk verwerkt! Dit hoofdstuk is eerder een
%% kort overzicht van je plan van aanpak.
%%
%% Maak voor elke fase (behalve het literatuuronderzoek) een NIEUW HOOFDSTUK aan
%% en geef het een gepaste titel.
\section{Research approach}

This research aims to develop and 
evaluate a solution for interpreting IBM MQ SMF data. 
The work is primarily design oriented, focusing on the creation of 
a classification framework and a supporting Proof of Concept (PoC).

Several research methods are applied throughout the different phases 
of the project:

\begin{itemize}
    \item \textbf{Literature study:} to understand IBM MQ buffer pool 
    behavior, SMF data structures, and existing performance indicators.
    
    \item \textbf{Requirement analysis:} to define relevant metrics and 
    expectations for the classification of operational states and the interpretation layer.
    
    \item \textbf{Data analysis:} to explore SMF 115 (subtype 215) data 
    and identify patterns, trends, and relevant indicators.
    
    \item \textbf{Design and implementation:} to develop the operational 
    state model and implement it in a Proof Of Concept.
    
    \item \textbf{Validation:} through manual comparison with MP1B output 
    and feedback from subject matter experts.
    
\end{itemize}






\section{Planning and deliverables}
\begin{table}[h]
    \centering
    \renewcommand{\arraystretch}{1.3}
    \begin{tabular}{|l|l|}
        \hline
        \textbf{Phase} & \textbf{Duration} \\
        \hline
         Problem Definition & 2 weeks  \\
        Data collection and \& Technical environment setup  & 2 weeks  \\
        State model design       & 2 weeks \\
        POC implementation          & 4 weeks \\
        Evaluation       & 2 weeks \\
        Reporting and reflection       & 2 weeks \\
        \hline    \end{tabular}
    \caption{Overview duration every phase}
\end{table}





\begin{table}[h]
    \centering
    \renewcommand{\arraystretch}{1.3}
    \begin{tabular}{|l|l|}
        \hline
        \textbf{Phase} & \textbf{Deliverable} \\
        \hline
        Problem Definition \& requirement analysis        &  buffer pool metrics\\
        &   requirement analysis  \\
        Data collection and \& Technical environment setup  & SMF 115, python, HTTP server, certificates  \\
        State model design       & Operational buffer pool  \\
        & state framework \\
        POC implementation          & Working Proof of concept \\
        Evaluation       & Evaluation report \\
        Reporting \& reflection      & Bachelor thesis document \\
        \hline
    \end{tabular}
    \caption{Overview deliverables every phase}
\end{table}







\section{Tools en technology}


\begin{table}[h]
    \centering
    \renewcommand{\arraystretch}{1.3}
    \begin{tabular}{|l|l|}
        \hline
        \textbf{Component} & \textbf{Tool} \\
        \hline
        Operating system & z/OS \\
        Platform      & IBM mainframe Z16 \\
        SMF extraction      & MP1B (MQSMF) that supports version IBM MQ v9.4 \\
        Job scheduler      & IBM Workload Scheduler  \\
        CSV transformation & Pandas v3.0.2 \\
        Analysis  & Pandas v3.0.2 \\
        Hosting  & HTTP Server \\
        Dataset format       & CSV \\
        Validation          & Manual MP1B analysis, SME validation \\
        Visualization       & Grafana v12.4.2 + Infinity plugin v3.8.0 \\
        \hline
    \end{tabular}
    \caption{Overview tools and technologies}
\end{table}

\subsection{SMF parser: MP1B (MQSMF)}

MP1B is an IBM supportpac that contains utility programs that are useful to test and work with IBM MQ objects. The supportpac contains 3 programs and is available to download at fixcentral. 
The first program is MQCMD and can be used to display queue and channel status' over time. The second program is OEMPUT and is used to get and put large amounts of messages to or from queues. OEMPUT is useful for throughput testing. The third program is MQSMF, which is used to parse SMF records into text and CSV files. 

This bachelor thesis uses the outputted CSV from the MQSMF program. The supportpac is at no cost with a valid IBM MQ license and offered as is. 

MP1B is an IBM Supportpac that contains an utility for extracting and formatting IBM MQ SMF statistics and accounting records. 
In this bachelor thesis MQSMF is used to convert SMF 115 (subtype 215) records into structured CSV datasets that are suitable for further processing. 
The tool was selected because it is offered at as is with an valid IBM MQ license. The supportpac is available to download at fixcentral.

Custom SMF parsers were not chosen due to it's irrelevance to the core of the thesis and PoC.

\subsection{Automated job scheduler: IBM Workload Scheduler}
Workload Scheduler is a job scheduling subsystem of z/OS that allows companies to automate and plan complex systems' workload. It schedules jobs according to requirements that can include time, dependencies, resources, events, calendars, priority, and many more. This POC makes use of an already existing scheduled job that executes the MQSMF utility program and writes the output to a mounted USS directory. 

\subsection{Analytical and data Manipulation tools: Pandas}
Pandas is a widely adopted Python library for data manipulation and numerical 
analysis. This PoC uses Pandas to group multiple datasets, transform raw CSV data, compute derived metrics, 
and implement the classification logic. Pandas was selected because of the CSV manipulation support, strong ecosystem and efficiency. 

\subsection{Hosting: IBM HTTP Server}
IBM HTTP server is used to host the processed data and a REST interface that has access to the data. This allows 
clients to retrieve the datasets through HTTPS protocol. The choice for IBM HTTP is due to the existing infrastructure and 
knowledge within the Rabobank, ensuring compatibility and ease of deployment.

\subsection{Data format: CSV}
CSV is used as the primary data format due to its compatibility with MP1B MQSMF output and Pandas. Additionally, CSV provides 
a lightweight and portable format that simplifies data exchange between components in  the PoC architecture.

JSON is also a valid format but would add complexity due to having to transform CSV to JSON standard. 
\subsection{Visualization: Grafana}
Grafana is used as the visualization layer of the PoC. It allows the presentation of statistics, operational states and trends in a clear way. 
Additionally Grafana chosen because of its flexibility, cost, and internal experience at Rabobank. The Infinity plugin is required to accept CSV formatted data.




